\documentclass{article}
\usepackage{amsmath}
\usepackage{amssymb}
\usepackage{amsthm}
\usepackage{hyperref}
\usepackage{url}

\newtheorem{theorem}{Theorem}
\newtheorem{proposition}{Proposition}
\newtheorem{remark}{Remark}

\title{Potential Shaping with Changing Coalitions:\\A Fixed-Population Ledger Repair}
\author{}
\date{}

\begin{document}
\maketitle

\begin{abstract}
MARS-RA uses multimodal pairwise comparisons, Bradley--Terry aggregation, and potential-based reward shaping for cooperative multi-agent reinforcement learning with changing active coalitions. Its comparison scores and potentials live in the current coalition space, so consecutive vector potentials need not have the same dimension; it also combines a scalar team reward with a vector shaping term without declaring payoff semantics. We give a narrow repair. Active credits are embedded in a fixed population space, inactive agents take null actions, and every population member retains a reward ledger on every transition. Componentwise telescoping then shifts each declared per-agent payoff by a policy-independent constant, preserving best responses and Nash equilibria for arbitrary comparison outputs. If shaping rewards are retained only while an agent is active, entry and exit boundary terms remain and may depend on policy. Nash invariance under static and time-varying potentials is prior work, and so is the observation that a nonzero potential at an action-dependent terminal state leaves a policy-dependent term; the contribution here is the typed ledger formulation for changing coalitions and the resulting entry and exit diagnosis. We also separate Bradley--Terry sampling error from semantic bias and explain why neither comparison accuracy nor a local signal bound yields a finite-training guarantee without a learner-stability premise.
% Ledger: 1, 3, 5, 7, 10, 12, 21, 24, 27, 28, 30, 31, 33. Prior-work attributions: search record, round 4.
\end{abstract}

\section{Introduction}

Learning-augmented algorithms require a declared prediction interface, an error measure, and a theorem connecting prediction quality to downstream performance. They distinguish consistency and error-sensitive guarantees from prediction-independent robustness, and warn that semantic predictions do not automatically compose into system guarantees \cite{zhao2026learning}.
% Ledger: 1, 2, 3, 13, 19.
MARS-RA instead uses multimodal pairwise comparisons to fit Bradley--Terry contribution scores and converts their softmax into a potential-based shaping signal for cooperative multi-agent reinforcement learning \cite{wang2026marsra}.
% Ledger: 1, 3, 9, 19.

The tempting question is how comparison error degrades the limiting return. That question is misplaced under exact potential shaping: the discounted shaping terms telescope, so their total is a common initial offset under the usual terminal and initial-potential conditions. The comparisons may affect optimisation, exploration, or sample efficiency, but comparison accuracy alone cannot control finite-training return without a stability, regret, or convergence contract for the learner.
% Ledger: 2, 4, 7, 9, 13, 15, 28, 31.

Our result is narrower. MARS-RA defines a fixed population but places each score in the current active-coalition space, then subtracts consecutive vector potentials. This expression is not typed when coalition sizes change. A fixed-population embedding is readily available, but inactive coordinates, transition ownership, and per-agent payoffs must be declared. We provide that construction and show exactly where an active-only implementation fails to telescope.
% Ledger: 10, 17, 21, 23, 24, 27, 30.

The main steps are not new. Lu, Schwartz, and Givigi prove that potential-based shaping preserves Nash equilibria in general-sum stochastic games \cite{lu2011policy}. MARS-RA itself appeals to a dynamic-potential extension by Devlin and Kudenko that covers potentials changing over time \cite{wang2026marsra}. Grze\'s showed that, in episodic tasks, a nonzero potential at an action-dependent terminal state leaves a term that can alter policies and Nash equilibria. Our contribution is the fixed-population reward-ledger formulation for changing coalitions, together with the observation that omitting inactive-agent ledger entries turns each agent's entry and exit into boundaries of the shaping sum, where terms of that kind reappear. We do not claim a new general reward-transformation theorem.
% Ledger basis for scope of the new claim: 21, 24, 25, 27, 28, 30, 31, 33. Prior-work attributions: search record, round 4.

\section{What the two papers establish}

The survey of Zhao et al. provides a vocabulary and discipline for prediction interfaces, robustness, error propagation, and prediction cost; it does not analyse MARS-RA or prove a MARL learner-stability theorem \cite{zhao2026learning}.
% Ledger: 2, 3, 4, 7, 13, 15, 19.
Wang et al. define MARS-RA over a fixed population with a changing active coalition. They obtain pairwise multimodal comparisons, aggregate them through a Bradley--Terry estimator, apply softmax, and use the resulting vector in potential-based reward shaping \cite{wang2026marsra}.
% Ledger: 1, 3, 10, 19, 21, 27, 30.

Their estimator statement concerns concentration around a latent Bradley--Terry preference \(c_t^*\), not a true contribution target \(v_t^*\). The identification \(c_t^*=v_t^*\) is an additional premise. Thus repeated queries can reduce sampling error while systematic semantic bias remains.
% Ledger: 1, 3, 5, 19.
Moreover, the supplied checks report that the stated estimator rate is not ready for composition: its population-size dependence is inconsistent across statements, connectedness alone does not supply uniform curvature, an unregularised maximum-likelihood estimate may diverge under separation, and the proof contains an approximate Taylor step. We therefore use no numerical Bradley--Terry rate.
% Ledger: 5, 8, 15, 25, 28, 31.

\section{Fixed-population formulation}

Let the population be \(N=\{1,\ldots,n\}\) and let \(I_t\subseteq N\) be the active coalition. Let \(\widehat c_t\) be any realised score vector indexed by \(I_t\). We permit biased, stochastic, or history-dependent scores. Any information needed to make the realised potential a state function may be included in an augmented state. This qualification is necessary for a standard stochastic-game equilibrium interpretation, although the algebra below is pathwise.
% Ledger: 7, 12, 24.

For nonterminal \(t\) with \(I_t\ne\varnothing\), define \(\Phi_t\in\mathbb R^n\) by
\begin{equation}
 \Phi_t^i=\mathbf 1\{i\in I_t\}
 \operatorname{softmax}_{I_t}(\widehat c_t)^i,
 \qquad i\in N,
 \label{eq:embedding}
\end{equation}
and set \(\Phi_T=0\). A complete model must separately declare a fallback for an empty coalition; nothing below supplies one. Give inactive agents null actions and maintain a payoff ledger for every \(i\in N\) on every transition. For declared per-agent environmental payoff \(r_t^i\), define the repaired team reward and objective by
\begin{equation}
 r_t^{\mathrm{team}}=\sum_{i\in N}r_t^i,
 \qquad
 J^{\mathrm{team}}(\pi)=
 \mathbb E_{\pi}\!\left[\sum_{t=0}^{T-1}\gamma^t
 r_t^{\mathrm{team}}\right].
 \label{eq:team-objective}
\end{equation}
If these payoffs allocate MARS-RA's original scalar team reward, this definition additionally requires their sum to equal that reward on each transition. Define
\begin{equation}
 F_t^i=\gamma\Phi_{t+1}^i-\Phi_t^i,
 \qquad \widetilde r_t^i=r_t^i+\rho F_t^i.
 \label{eq:shaping}
\end{equation}
% Ledger: 24, 27, 30.

\begin{theorem}[fixed-population ledger invariance]
For a fixed initial augmented state and terminal potential \(\Phi_T=0\), the construction in \eqref{eq:embedding}--\eqref{eq:shaping} shifts player \(i\)'s discounted payoff by \(-\rho\Phi_0^i\). Consequently, it preserves every unilateral payoff difference, best-response correspondence, and Nash equilibrium. It also preserves the ordering of joint policies under the repaired team objective \(J^{\mathrm{team}}\) in \eqref{eq:team-objective}. These conclusions hold for arbitrary realised comparison outputs on a common pathwise realisation.
% Ledger: 24, 27, 33.
\end{theorem}

\begin{proof}
For each \(i\in N\),
\begin{align*}
 \sum_{t=0}^{T-1}\gamma^t F_t^i
 &=\sum_{t=0}^{T-1}
   \left(\gamma^{t+1}\Phi_{t+1}^i-\gamma^t\Phi_t^i\right)\\
 &=-\Phi_0^i+\gamma^T\Phi_T^i=-\Phi_0^i.
\end{align*}
The initial augmented state is fixed, so the final quantity is independent of every player's policy. Adding a constant to player \(i\)'s payoff preserves all unilateral payoff differences and hence its best responses. Applying this to every player preserves Nash equilibria. Summing the componentwise identities gives the team-policy statement.
% Ledger: 4, 7, 12, 24, 27.
\end{proof}

This theorem is a proposed repair, not a theorem stated by MARS-RA. Its equilibrium-invariance logic is an instance of the established potential-shaping result for stochastic games \cite{lu2011policy}, and its allowance for potentials that change over time matches the dynamic-shaping extension that MARS-RA cites \cite{wang2026marsra}. What is specific here is the fixed index space and the all-population transition ledger.
% Ledger: 21, 24, 27, 28, 31. Prior-work attributions: search record, round 4.

\section{The active-only boundary}

Suppose instead that player \(i\) retains shaping payments only during an active interval \([a,b)\). Direct summation gives boundary terms of the form
\begin{equation}
 \sum_{t=a}^{b-1}\gamma^t F_t^i
 =-\gamma^a\Phi_a^i+\gamma^b\Phi_b^i,
 \label{eq:boundary}
\end{equation}
with the second term adjusted according to ownership of the exit transition. Entry, exit, and coalition changes can depend on earlier actions, so \eqref{eq:boundary} need not be policy-independent. Active-only equivalence therefore requires explicit transition ownership or compensating transfers that reproduce the fixed ledger sum. The inputs do not specify those transfers, so we leave their design open.
% Ledger: 24, 25, 27, 28, 31.
The exit term plays, for one agent's participation interval, the role of the terminal-potential term that Grze\'s analysed for episodic tasks. The entry term has no counterpart in that analysis, where the initial potential does not depend on actions.
% Comparison with published work: search record, round 4.

There is a second interface boundary. MARS-RA begins from a scalar team reward but adds a vector shaping term. One coherent repair declares per-agent payoffs and applies the componentwise theorem. Another uses a scalar team potential, which preserves the ordering of joint policies under the team objective. The inputs examine only one such potential, the coordinate sum of the embedded credits. Since \(\sum_i\Phi_t^i=1\) for every nonempty nonterminal coalition, coordinate-sum scalarisation yields a potential that does not depend on the comparisons, so its shaping signal carries no comparison information.
% Ledger: 30, 33.
This observation is narrower than its wording in the supporting notes, which conclude from it that comparisons can matter only through redistributed per-agent signals. The argument supports that conclusion only against coordinate-sum scalarisation. Other scalar functions of the credit vector, such as a weighted sum with unequal weights, are in general not constant, can depend on the comparisons, and still define a scalar potential. The notes do not analyse them, and we do not exclude comparison-dependent scalar team potentials.
% The over-broad wording is in ledger entries 30 and 33 and in the note Q1P6.tex. Scope restriction following review round 3, finding F1; not a result of the inputs.

\section{Estimation, semantics, and local signal sensitivity}

After fixing the translation gauge of Bradley--Terry scores, define semantic bias
\begin{equation}
 \delta_t=\lVert c_t^*-v_t^*\rVert_2.
\end{equation}
Then
\begin{equation}
 \lVert\widehat c_t-v_t^*\rVert_2
 \leq \lVert\widehat c_t-c_t^*\rVert_2+\delta_t.
 \label{eq:bias}
\end{equation}
Equation \eqref{eq:bias} is only the triangle inequality. It separates sampling error about the comparator's latent preference from semantic error to the declared contribution target; it does not validate that target or repair the estimator theorem.
% Ledger: 1, 3, 5, 8, 19.

For this local comparison, assume that the estimated and target scores at each time are indexed by the same nonempty active coalition \(I_t\). Let
\begin{equation}
 d_t=\widehat c_t-v_t^*\in\mathbb R^{|I_t|},
\end{equation}
after fixing a common translation gauge. When it appears in the shaping-signal bound below, regard \(d_t\) as zero-coordinate embedded in \(\mathbb R^n\), and set \(d_T=0\in\mathbb R^n\) by convention. Embed the target credits in the same fixed population space by
\begin{equation}
 (\Phi_t^*)^i=\mathbf 1\{i\in I_t\}
 \operatorname{softmax}_{I_t}(v_t^*)^i,
 \qquad \Phi_T^*=0,
\end{equation}
and define the vector shaping signals
\begin{equation}
 \widehat F_t=\gamma\Phi_{t+1}-\Phi_t,
 \qquad F_t^*=\gamma\Phi_{t+1}^*-\Phi_t^*.
\end{equation}
The zero-coordinate embeddings preserve Euclidean distances on their common active coalition. Since softmax is \(1/2\)-Lipschitz in Euclidean norm, these one-step shaping signals satisfy
\begin{equation}
 \lVert\widehat F_t-F_t^*\rVert_2
 \leq \frac12\left(
  \gamma\lVert d_{t+1}\rVert_2+\lVert d_t\rVert_2
 \right).
 \label{eq:local}
\end{equation}
This is a local reward-signal bound only. It does not bound the return of a policy produced by MAPPO after finitely many updates. Learners can react differently to the same shaped signal, and neither source paper provides the missing stability, regret, or convergence premise.
% Ledger: 5, 7, 9, 13, 15, 28, 31.

\section{Related work}

Lu, Schwartz, and Givigi extend potential-based shaping to multi-player general-sum stochastic games and prove preservation of Nash equilibria \cite{lu2011policy}. Their paper also reports similar contemporaneous work by Devlin and Kudenko. Accordingly, the telescoping and general Nash-invariance argument is already published and is not a contribution of this paper.

Two further published results cover steps used here. Devlin and Kudenko (AAMAS 2012) let the potential depend on the time at which a state is visited and show, by the same telescoping along a fixed sequence, that the shaped return differs from the unshaped return only by the initial potential; policy invariance and consistent Nash equilibria therefore survive time-varying potentials. MARS-RA cites this result in support of its own shaping \cite{wang2026marsra}. It covers the clause of our theorem that allows arbitrary realised comparison outputs. Grze\'s (AAMAS 2017) shows that in episodic tasks the shaped return also contains the discounted potential of the state at which the trajectory stops, that this term depends on actions when termination does, and that unless this potential is zero it can alter the optimal policy and, in a two-agent coordination game, the set of Nash equilibria. The exit term in \eqref{eq:boundary} is a term of this kind for a single agent's participation interval, so the main step of our boundary diagnosis is also published; what is added is its placement at the entry and exit of individual agents under changing coalitions. We read both papers in their author-deposited versions. We name them in the text rather than in the bibliography because we could not verify their bibliographic records against a DOI landing page or an arXiv abstract page.

Our literature search did not locate a verified publication that states the particular fixed-population, inactive-coordinate, all-transition ledger construction for changing active coalitions, or that treats individual entry and exit as boundaries of the shaping sum. This is a search result, not a novelty proof.
% Ledger basis for delimiting the supported pair-specific result: 21, 24, 25, 27, 28, 30, 31, 33. Prior-work attributions: search record, rounds 1 to 4.

The two focal papers supply the remaining context. Zhao et al. insist on typed semantic interfaces and explicit downstream composition \cite{zhao2026learning}; Wang et al. supply the changing-coalition comparison and shaping mechanism to which the interface audit applies \cite{wang2026marsra}.
% Ledger: 1, 3, 10, 19, 21.

\section{Interpretation}

The repair makes objective-level robustness prediction-independent: once payoff semantics and ledger ownership are fixed, even inaccurate comparisons preserve payoff differences. This should not be read as evidence that comparisons help learning. With a fixed initial augmented state and zero terminal potential, exact shaping changes discounted return only by an initial offset, whether the potential is per-agent or scalar, so any benefit lies on the finite-training optimisation path. Under the per-agent repair that path is driven by redistributed per-agent signals. We do not claim that this is the only route: the coordinate-sum observation does not rule out a comparison-dependent scalar team potential, and the inputs do not examine its effect on training. Better estimation can reduce one component of signal error while sharpening a semantically biased target.
% Ledger: 1, 3, 4, 5, 7, 13, 28, 30, 31, 33.

\section{Limitations and open questions}

The result is algebraic and conditional on declared per-agent payoffs, a fixed initial augmented state, a terminal zero potential, null inactive actions, and an all-population transition ledger. It is undefined when \(I_t=\varnothing\) unless the model adds an explicit empty-coalition potential convention. It does not establish that the published MARS-RA implementation uses this ledger. It does not give compensating entry or exit transfers for active-only storage. For stochastic or history-dependent comparison outputs, an equilibrium theorem may require state augmentation beyond the pathwise identity. The scalarisation observation covers only the coordinate sum of the embedded credits; comparison-dependent scalar team potentials are neither excluded nor analysed. Two of the prior works that cover our main steps are named in the text but could not be entered in the bibliography.
% Ledger: 12, 24, 25, 27, 28, 30, 31, 33.

We do not provide a corrected Bradley--Terry concentration rate, a MAPPO stability theorem, a finite-training return bound, or a query-cost trade-off. These remain open, as do common units for pricing multimodal-model queries against return or sample efficiency, and whether a comparison-dependent scalar team potential gives a useful training signal. The fixed-space theorem is a correctness layer, not evidence of a positive performance gain and not a general theory of open Markov games.
% Ledger: 8, 15, 25, 28, 31.

\section{Conclusion}

Changing coalitions expose a typing and accounting requirement hidden by the usual potential-shaping notation. A fixed-population all-transition ledger restores componentwise telescoping, by the published argument applied in a fixed index space, and preserves payoff differences for arbitrary comparison outputs. Active-only accounting leaves entry and exit boundary terms of the kind already known for episodic terminal states. Scalarising the credits by their coordinate sum removes the comparison signal; other scalarisations are not analysed here. Under exact shaping, any benefit of comparison-based shaping lies in finite-training dynamics, for which a learner contract is still missing.
% Ledger: 4, 7, 24, 25, 27, 28, 30, 31, 33. Prior-work attributions: search record, round 4.

\bibliographystyle{plain}
\bibliography{references}

\end{document}
